\documentclass[11pt,a4paper,oneside,openany]{ltjsbook}
\input{preamble.tex}
\geometry{top=25mm,bottom=26mm,inner=30mm,outer=25mm}
\hypersetup{pdftitle={Scientific Monograph (demonstration)},pdfauthor={Research Team}}
\fancyhead[L]{\footnotesize\sffamily\color{researchgray}\nouppercase{\leftmark}}
\begin{document}
\begin{titlepage}
\centering
\vspace*{42mm}
{\sffamily\large\color{researchteal} RESEARCH MONOGRAPH \par}
\vspace{12mm}
{\sffamily\bfseries\fontsize{25}{33}\selectfont\color{researchnavy}表現学習の基礎と\\[4mm]実験的評価\par}
\vspace{12mm}
{\Large Foundations and Evaluation\par}
\vfill
{\large Research Team \quad Laboratory / Project\par}
\vspace{8mm}
{\large 2026年10月15日\par}
\vspace{5mm}
{\color{researchgray}\small 本文と数値は組版確認のための例示であり、研究成果ではありません。\par}
\end{titlepage}
\frontmatter
\chapter*{はじめに}
本書は技術要覧や研究チュートリアルのためのテンプレートである。
表示例に含まれる学習曲線と数値は全て\textbf{架空データ}である。
\tableofcontents
\mainmatter
\chapter{背景と定義}
\section{表現学習の目的}
データ $x$ を特徴表現 $z=f_\theta(x)$ に写像することを考える。
記号の意味、仮定、対象データの適用範囲を冒頭で明記する。
化学表記法の参考文献としてSMILESの原著\cite{weininger1988smiles}を示す。
\begin{definition}[表現写像]
入力空間 $\mathcal{X}$ から潜在空間 $\mathcal{Z}$ への写像 $f_\theta:\mathcal{X}\to\mathcal{Z}$ を表現写像という。
\end{definition}
\begin{researchnote}[title={読み進める際の注意}]
本書に示すグラフは、学術的な事実の主張ではなく組版機能のデモンストレーションである。
\end{researchnote}
\section{損失関数}
系列 $x_{1:T}$ の平均的な負の対数尤度を次のように定義する。
\begin{equation}\label{eq:loss}
\mathcal{L}(\theta)=-\frac{1}{T}\sum_{t=1}^T\log p_\theta(x_t\mid x_{<t}).
\end{equation}
\chapter{実験的評価の設計}
\section{比較条件とグラフ}
\begin{figure}[htbp]
\centering
\includegraphics[width=0.84\linewidth]{figures/example-curve.pdf}
\caption{表示例として作成した合成データの損失曲線。}\label{fig:curve}
\end{figure}
\cref{fig:curve}でベクター図の品質を確認する。実運用では、評価データと作図スクリプトを両方保存する。
\section{仮説と検証の分離}
観測結果は数値と一緒に記し、解釈・仮説・今後の検証は別の段落に記す。
\begin{table}[htbp]
\centering
\caption{条件別の架空損失値（検証対象ではない）}\label{tab:results}
\begin{tabular}{lSS}
\toprule
{条件} & {初期値} & {最終値}\\
\midrule
Baseline & 1.75 & 0.73\\
Variant A & 1.68 & 0.49\\
Variant B & 1.60 & 0.35\\
\bottomrule
\end{tabular}
\end{table}
\chapter{再現性と記録}
\section{必須記録}
\begin{enumerate}
\item データセット識別子と版、学習・評価分割。
\item コードのコミットID、環境、ハイパーパラメータ。
\item 評価指標の定義、シード、統計の扱い。
\end{enumerate}
\backmatter
\printbibliography[title={参考文献},heading=bibintoc]
\end{document}
